\clearpage
\phantomsection
\section*{Chương 8. Chuẩn hóa $C[\tau]=\mathcal{R}_w[0]/\mathcal{R}_w[\tau]$ (ROM nhân thay chia)}
\addcontentsline{toc}{section}{Chương 8. Chuẩn hóa C[τ] (ROM nhân thay chia)}
\pdfbookmark[1]{Chương 8. Chuẩn hóa C[τ] (ROM nhân thay chia)}{chuong8}

\begin{learningobj}
Sau chương này bạn giải thích được vì sao phải chia cho tự tương quan của cửa sổ; dẫn ra $C[\tau]$ từ nguyên lý đầu tiên; tính tay và kiểm Python một bảng nhỏ cùng các mốc trên khung $L=640$; cấu hình ROM M10K mang $C[\tau]$ dạng Q3.13; đối chiếu với LPM\_DIVIDE / bộ chia non-restoring khi $\mathcal{R}_w$ đổi lúc chạy; viết twin sinh \texttt{.hex} bit-exact; và đọc được hộp chi phí --- tại sao nhân thay chia runtime.
\end{learningobj}

\subsection*{8.1. Hiện tượng thật: đỉnh $\tau$ lớn bị ``đạp'' vì hai quả núi ít chồng}

Sau khi nhân cửa sổ Hanning (Chương~6), mỗi khung trở thành một ``quả núi'': biên độ lớn ở giữa, nhỏ ở hai đầu. Khi tính $\mathcal{R}_x[\tau]$, ta đang đo độ chồng giữa quả núi đó và bản sao trễ $\tau$ mẫu. Trễ càng lớn, hai quả núi chồng nhau càng ít --- năng lượng chồng đó chính là tự tương quan của \emph{chính cửa sổ}:
\[
\mathcal{R}_w[\tau] = \sum_{n=0}^{L-1-\tau} w[n]\, w[n+\tau].
\]
Nếu lấy $\mathcal{R}_x[\tau]$ thô làm điểm số, đỉnh ở $\tau$ lớn (giọng trầm, $T_0$ dài) bị đạp xuống chỉ vì hình học cửa sổ, không phải vì tín hiệu kém tuần hoàn. Ví dụ khung cosin ở Chương~7 cho $\mathcal{R}[8]/\mathcal{R}[0] = 0{,}5$ dù sóng lặp hoàn hảo --- đó đúng là hiệu ứng này.

Praat/Boersma nhân hệ số bù
\[
C[\tau] = \frac{\mathcal{R}_w[0]}{\mathcal{R}_w[\tau]}
\]
để kéo mọi đỉnh về cùng thước gần $[0,1]$. \textbf{$C[\tau]$ không phải hiện tượng dây thanh}: nó là hiệu chỉnh hình học của cửa sổ. Trên FPGA, $\mathcal{R}_w$ chỉ phụ thuộc $w$ và $L$ (bất biến lúc chạy), nên tính offline, nạp ROM, rồi \emph{nhân} thay vì chia mỗi $\tau$.

\begin{physicalinsight}
Quên nhân $C[\tau]$ thiên lệch hệ thống về $\tau$ nhỏ (giọng cao giả): đỉnh trầm bị đạp, thuật toán thích đỉnh ngắn hơn. Đây là lỗi kiến trúc, không phải nhiễu ngẫu nhiên.
\end{physicalinsight}

\begin{figure}[H]
\centering
\begin{tikzpicture}[font=\small, >=Stealth]
  % two overlapping mountains
  \draw[thick, teal, opacity=0.85] (0.5,0) .. controls (1.5,2.4) and (2.5,2.4) .. (3.5,0);
  \draw[thick, orange!80!black, opacity=0.85] (2.0,0) .. controls (3.0,2.4) and (4.0,2.4) .. (5.0,0);
  \draw[<->, gray] (1.8,0.15)--(3.2,0.15);
  \node[gray, font=\scriptsize] at (2.5,-0.35){vùng chồng $\propto \mathcal{R}_w[\tau]$};
  \node[teal, font=\scriptsize] at (1.5,2.6){$w[n]$};
  \node[orange!80!black, font=\scriptsize] at (4.0,2.6){$w[n+\tau]$};
  \draw[->] (5.8,1.2)--(7.0,1.2);
  \node[font=\scriptsize, align=left, anchor=west] at (7.1,1.5){$C[\tau]=\dfrac{\mathcal{R}_w[0]}{\mathcal{R}_w[\tau]}$\\nâng đỉnh $\tau$ lớn};
\end{tikzpicture}
\caption{Hai quả núi cửa sổ chồng ít dần khi $\tau$ tăng. Hệ số $C[\tau]$ bù đúng độ mất chồng đó.}
\label{fig:ch08_window_overlap}
\end{figure}

\subsection*{8.2. Toán từng bước --- mọi ký hiệu}

Công thức chuẩn hóa Boersma (sau khi đã có $\mathcal{R}_x[0]$):
\[
r_a[\tau] = \frac{\mathcal{R}_x[\tau]}{\mathcal{R}_x[0]}\cdot C[\tau], \qquad
C[\tau] = \frac{\mathcal{R}_w[0]}{\mathcal{R}_w[\tau]}.
\]
\begin{itemize}[leftmargin=1.4em]
\item $\mathcal{R}_x[\tau]$: tự tương quan tín hiệu đã cửa sổ (Chương~7), int rộng trên silicon.
\item $\mathcal{R}_w[\tau]$: tự tương quan cửa sổ --- mẫu số ``độ chồng'', tính offline từ $w_q$.
\item $C[\tau]\ge 1$ và tăng theo $\tau$ (vì $\mathcal{R}_w$ giảm).
\item Lượng tử silicon: $C_q[\tau] = \mathrm{sat}_{u16}\!\bigl(\mathrm{round}(C[\tau]\cdot 2^{13})\bigr)$ --- định dạng \textbf{Q3.13} (3 bit phần nguyên gồm đầu, 13 bit thập phân; $1{,}0\mapsto 8192$; dải đến gần $4{,}0$). Chọn Q3.13 vì trên $L=640$, $C[214]\approx 2{,}13$ vượt quá 1 (Q1.15 sẽ bão hòa sai).
\item Bit-growth: $R_x$ (vd Q2.30 / accum) $\times C$ (Q3.13) rồi làm tròn về thang điểm số. Twin và RTL dùng cùng quy tắc:
\[
S[\tau] = (R_x[\tau]\cdot C_q[\tau] + 2^{12}) \gg 13
\]
(làm tròn gần nhất, khớp \texttt{lab\_1\_3\_5.py}).
\end{itemize}

Dải ROM phần cứng: 189 từ cho $\tau = 26\ldots 214$. Tệp \texttt{hex/c\_tau\_rom.hex} trong gói này còn nhân thêm hệ số octave-cost của Praat (xem \texttt{scoring\_rom\_q13} trong golden) --- đó là tiện ích gộp ROM, không đổi bản chất $C[\tau]$ bù cửa sổ. Twin kiểm tra cửa sổ thuần dùng \texttt{compensation\_rom\_q13}.

\subsection*{8.3. Ví dụ tính tay --- mọi số đã chạy Python}

\paragraph{A. $L=4$, $w=[1,2,2,1]$.}
\begin{align*}
\mathcal{R}_w[0] &= 1+4+4+1 = 10,\\
\mathcal{R}_w[1] &= 1\cdot 2 + 2\cdot 2 + 2\cdot 1 = 8,\\
\mathcal{R}_w[2] &= 1\cdot 2 + 2\cdot 1 = 4,\\
\mathcal{R}_w[3] &= 1\cdot 1 = 1.
\end{align*}
Vậy $C = [1,\ 1{,}25,\ 2{,}5,\ 10]$. Q3.13: $1{,}25 \mapsto \mathrm{round}(1{,}25\cdot 8192) = 10240$. Giả sử $\mathcal{R}_x = [20,12,6,1]$ thì sau nhân $C$: $[20,\ 15,\ 15,\ 10]$ --- đỉnh phụ $\tau$ được nâng tương đối.

\paragraph{B. $w=[1,2,1]$, $\tau=2$.}
$\mathcal{R}_w[0] = 6$, $\mathcal{R}_w[2] = 1$, $C[2] = 6$.

\paragraph{C. Khung thật $N=640$ Praat (bù cửa sổ thuần).}
\begin{center}\small
\begin{tabular}{@{}rrrl@{}}\toprule
$\tau$ & $C$ float & $C_q$ Q3.13 & hex \\
\midrule
26 & 1.010885 & 8281 & \texttt{2059} \\
80 & 1.108152 & 9078 & \texttt{2376} \\
160 & 1.515087 & 12412 & \texttt{307C} \\
214 & 2.127577 & 17429 & \texttt{4415} \\
\bottomrule
\end{tabular}
\end{center}
Hex triển khai (\texttt{scoring\_rom}, có octave-cost): $C_q[26] = 8283$, $C_q[214] = 17070$ --- lệch nhẹ so với bảng thuần, đúng thiết kế gộp phạt quãng tám vào cùng một lần nhân.

\begin{figure}[H]
\centering
\begin{tikzpicture}[box/.style={draw,align=center,font=\small,rounded corners,minimum height=9mm}, node distance=7mm]
\node[box] (mac) {MAC $\mathcal{R}_x[\tau]$};
\node[box, right=of mac] (rom) {ROM $C[\tau]$\\Q3.13};
\node[box, right=of rom] (mul) {DSP $\times$};
\node[box, right=of mul] (peak) {argmax / parabolic};
\draw[->,thick] (mac)--(mul); \draw[->,thick] (rom)--(mul); \draw[->,thick] (mul)--(peak);
\end{tikzpicture}
\caption{Sau nhân $C$: đưa $S[\tau]$ (hoặc $r_a[\tau]$) vào khối dò đỉnh (Chương~9).}
\label{fig:ch08_norm_datapath}
\end{figure}

\subsection*{8.4. Ý nghĩa phần cứng --- ROM mặc định, chia chỉ khi bắt buộc}

Đường mặc định: IP \textbf{On-Chip Memory} kiểu ROM (hoặc RAM khởi tạo \texttt{.hex}/\texttt{.mif}).
\begin{itemize}[leftmargin=1.4em]
\item Depth $= 189$ (hoặc $\tau_{\max}+1$ nếu lưu từ 0); width $= 16$ bit mang Q3.13.
\item Registered output: 1 chu kỳ --- FSM Chương~11 chờ đúng latency trước khi nhân.
\item Một khối DSP nhân $R_x\times C_q$ mỗi $\tau$ (có thể dùng chung nhịp với MAC).
\end{itemize}

Khi nào mới chia runtime? Chỉ khi $\mathcal{R}_w$ \emph{đổi lúc chạy} (ví dụ đổi loại cửa sổ động). Khi đó dùng IP \textbf{LPM\_DIVIDE} (User Guide chia số nguyên của Intel) hoặc bộ chia \textbf{non-restoring} tự viết bằng ALM --- cùng họ chia sẽ gặp lại ở Chương~9 cho $F_0 = f_s/\tau$. Đường pitch chuẩn của tài liệu này \textbf{không} chia mỗi $\tau$: ROM đã nướng sẵn nghịch đảo.

Đổi $L$ hoặc công thức $w$ thì \textbf{rebuild file hex bắt buộc}; nạp ROM cũ sẽ sai im lặng và đỉnh $F_0$ thiên lệch.

\paragraph{RTL (ý tưởng).}
\begin{verbatim}
(* ram_init_file = "c_tau_rom.hex" *)
reg [15:0] C_rom [0:188];           // tau 26..214
reg [15:0] C_q;
reg signed [42:0] S;                // score after *C
always @(posedge clk) begin
  C_q <= C_rom[tau - 26];           // 1-cycle registered read
  S   <= (R_x * $signed({1'b0,C_q}) + 43'sd4096) >>> 13;
end
\end{verbatim}

\subsection*{8.5. Twin Python bit-exact}

\begin{verbatim}
import numpy as np
Q13 = 1 << 13

def build_C(w, tau_min, tau_max):
    w = np.asarray(w, float)
    L = len(w)
    Rw0 = float(np.dot(w, w))
    out = []
    for t in range(tau_min, tau_max + 1):
        Rwt = float(np.dot(w[: L - t], w[t:]))
        out.append(int(round(Rw0 / max(Rwt, 1e-30) * Q13)))
    return np.clip(out, 0, 65535).astype(np.uint16)

def apply_C(R_x, C_q):
    return (R_x.astype(np.int64) * C_q.astype(np.int64) + (1 << 12)) >> 13

# assert: build_C([1,2,2,1], 0, 3)[1] == 10240  # C=1.25
\end{verbatim}

\subsection*{8.6. Chi phí silicon}

\begin{costbox}
\textbf{Ngân sách chuẩn hóa $C[\tau]$.}
\begin{itemize}[leftmargin=1.2em]
\item \textbf{M10K}: $189\times 16 = 3\,024$\,bit $< 1$ M10K ($10\,240$\,bit).
\item \textbf{DSP}: 1 nhân $R_x\times C$ mỗi $\tau$ (dùng chung hoặc ngay sau MAC).
\item \textbf{ALM}: địa chỉ + đăng ký --- vài chục.
\item So với LPM\_DIVIDE mỗi $\tau$: tiết kiệm 189 phép chia/khung và latency không đều --- lý do chọn ROM.
\item Vì sao Q3.13: $C[214]\approx 2{,}13$ vừa trong 3 bit nguyên; Q1.15 sẽ bão hòa tại $1{,}0$ và phá bù.
\end{itemize}
\end{costbox}

\subsection*{8.7. Bài tập và lời giải giải thích}

\paragraph{Bài tập 8.1.} $w=[1,2,1]$: tính $C[2]$.

\paragraph{Bài tập 8.2.} Hex Q3.13 cho $C = 1{,}5$: giá trị nguyên và mã hex?

\paragraph{Bài tập 8.3.} Twin dùng floor thay round khi lượng tử $C$: lệch tối đa bao nhiêu LSB trước khi cắt?

\paragraph{Bài tập 8.4.} Quên nhân $C$: đỉnh $\tau$ lớn bị thiên lệch thế nào?

\paragraph{Bài tập 8.5.} Liệt kê ba tham số IP ROM bắt buộc: width, depth, init file. Khi nào cân nhắc LPM\_DIVIDE thay ROM?

\noindent\textbf{Lời giải giải thích.}
\begin{enumerate}[leftmargin=1.4em]
\item $\mathcal{R}_w[0] = 1+4+1 = 6$, $\mathcal{R}_w[2] = 1\cdot 1 = 1$, $C[2] = 6$.
\item $1{,}5\cdot 8192 = 12288 = $\texttt{0x3000}.
\item Floor thay round lệch tối đa $0{,}5$ LSB của thang Q3.13 trước khi ép kiểu nguyên.
\item Quên $C$: đỉnh $\tau$ lớn bị đạp bởi hình học cửa sổ, thuật toán thiên về $\tau$ nhỏ (giọng cao giả).
\item Width 16, depth 189 (hoặc $\tau_{\max}-\tau_{\min}+1$), init \texttt{c\_tau\_rom.hex}. LPM\_DIVIDE chỉ khi $\mathcal{R}_w$ đổi lúc chạy; đường mặc định pitch không chia mỗi $\tau$.
\end{enumerate}
